\section{Conclusion}
\label{sec:conclusion}

We fixed a leakage-controlled evaluation protocol, held it constant,
and varied individual evaluation choices while keeping the model, data
and training procedure identical: five architectures, three
technologies, three horizons, two regimes and five seeds, 900 runs
across a validation and a held-out test year.

The protocol effects are one to two orders of magnitude larger than the
architectural ones. Reference forecast: 70 percent of reported skill at
six hours. Night-hour inclusion: approximately 34 percent deflation in
normalised error. Residual-stage construction: 0.034 becomes an
apparent 0.334. Architecture: the largest difference we measure between
any two architectures is 0.024.

The architectural findings are correspondingly modest. Gradient
boosting and a recurrent network are statistically indistinguishable at
one- and three-hour horizons. A convolutional front end shows no
detectable benefit. A gradient-boosted residual correction stage
reduces skill in all eighteen configurations tested.

Every finding holds in direction on the held-out test year, evaluated
once, after all choices were frozen. Magnitudes shift, most notably a
denominator effect that moves two co-located arrays' skill in opposite
directions between years, because their optimal reference weights
differ.

\subsection{Implications for Practice}
\label{sec:conclusion-implications}

Three recommendations follow, none novel, most absent from the
literature we surveyed. Report a skill score against a defensible
reference and state which reference: the two references in this paper
support opposite conclusions from identical forecasts. Report the
evaluation sample: night-hour inclusion, the daylight threshold, and
excluded hours each move reported error by more than the improvement
typically claimed. Report run-to-run variance: none of the 27 papers we
surveyed does, and the seed spread we measure (0.017) is of the same
order as the largest architecture difference.

We add one recommendation specific to multi-stage architectures: state
where residual-stage training data come from. Fitting on the evaluation
split is in-sample evaluation of that stage, and in our measurements it
reverses the sign of the component's contribution.

\subsection{Scope}
\label{sec:conclusion-scope}

These results come from three co-located arrays at a single desert
site, sharing one weather station, over one held-out year, at horizons
of one to six hours, with deterministic forecasts and no hyperparameter
tuning. We do not claim that any architecture is generally inferior,
nor that these effect sizes transfer to other sites or climates. We
claim that the choices producing them are made in every study of this
kind, are unreported in most, and are large enough that reported
accuracies cannot be compared across papers without them.

Rigorous verification practice for solar forecasting exists and is
documented \cite{yang2020verification, yang2019ropes}. One of the 27
papers we surveyed follows it \cite{mayer2022physmlhybrid}. The
harness, the run records, the coded survey with its supporting
evidence, and the scripts regenerating every table and figure in this
paper are available at an anonymised repository, link withheld for
review, so that the protocol described here can be adopted, or
disputed, on the same data.
